\subsection{Recovering the observations and the spin factor}
\label{sec:recovery}

The comparison has now proved $V(t)=0$ for the particular copies
sampled in~\eqref{eq:copy-law}: Proposition~\ref{prop:strict-uniqueness}
handles $b\theta^2<1$, and Proposition~\ref{prop:critical-uniqueness}
handles $b\theta^2=1$. We first deduce equality of their observation
paths and turn it into a Borel map from innovations to observations.
The spins will then be read from the integer-time slopes of those paths.

The final passage to a spin map equivariant on every input uses local
conditional expectations.

\begin{lemma}[Pointwise equivariance of spin factors]
\label{lem:pointwise-spin-factor}
Let $(U_v)_{v\in\mathcal V}$ be independent uniform labels on $T_d$.
Suppose a Borel map
$\Phi:[0,1]^{\mathcal V}\to\{-1,1\}^{\mathcal V}$ satisfies
$\Phi(gU)=g\Phi(U)$ almost surely for each fixed
$g\in\Aut(T_d)$. There is a Borel map
$\Psi:[0,1]^{\mathcal V}\to\{-1,1\}^{\mathcal V}$ such that
$\Psi(U)=\Phi(U)$ almost surely and
\[
 \Psi(gx)=g\Psi(x)
 \qquad\text{for every }x\in[0,1]^{\mathcal V}
                 \text{ and }g\in\Aut(T_d).
\]
\end{lemma}

\begin{proof}
Fix an auxiliary vertex $o$. For each integer $R\ge0$, choose a
Borel function $h_R:[0,1]^{B_R(o)}\to[-1,1]$ representing
\[
 \E[\Phi(U)_o\mid U|_{B_R(o)}].
\]
A Borel version exists, and clipping it to $[-1,1]$ preserves this
conditional expectation. Let $G_R$ be the finite automorphism group
of the rooted ball $B_R(o)$, and average:
\[
 q_R(y)=\frac1{|G_R|}\sum_{a\in G_R}h_R(ay).
\]
Every $a\in G_R$ extends to a tree automorphism fixing $o$: outside
the ball, match the isomorphic rooted components attached to its
boundary. Such an extension preserves the IID law and the ball's
label sigma-field. The assumed almost-sure equivariance therefore
makes $h_R(a(U|_{B_R(o)}))$ another version of the same conditional
expectation. Thus $q_R$ still represents that conditional expectation,
and is Borel, bounded by one, and invariant under $G_R$ on every input.

Apply this same local rule at every vertex. Explicitly, for a rooted
isomorphism $\iota:B_R(o)\to B_R(v)$, put
\[
 A_R(x)_v=q_R\bigl((x_{\iota(u)})_{u\in B_R(o)}\bigr).
\]
Two choices of $\iota$ differ by an element of $G_R$, so the value
does not depend on this choice. The maps $A_R$ are consequently
Borel and satisfy $A_R(gx)=gA_R(x)$ for every input and automorphism.
Each $\iota$ also extends to a tree automorphism taking $o$ to $v$.
Transporting the conditional expectation by this fixed automorphism
shows that
\[
 A_R(U)_v=\E[\Phi(U)_v\mid U|_{B_R(v)}]
                  \quad\text{almost surely}.
\]
The ball sigma-fields increase to the full label sigma-field.
Upward martingale convergence, together with a countable intersection
over vertices and radii, gives $A_R(U)_v\to\Phi(U)_v$ simultaneously
at every vertex almost surely.

On every label input define
\[
 \Psi(x)_v=
 \begin{cases}
  1,&\limsup_{R\to\infty} A_R(x)_v\ge0,\\
 -1,&\limsup_{R\to\infty} A_R(x)_v<0.
 \end{cases}
\]
Boundedness makes each limsup finite, and this formula defines a
Borel spin map. The pointwise equivariance of every $A_R$ passes
through the limsup and the common threshold. On the simultaneous
convergence event the limits belong to $\{-1,1\}$, so
$\Psi(U)=\Phi(U)$ almost surely.
\end{proof}

\begin{proposition}[Recovery of observations and spins]
\label{prop:factor-extraction}
Fix $d\ge3$ and $0<\theta=\tanh\beta\le(d-1)^{-1/2}$, and retain
the path spaces and laws $\mathsf Q,\mathsf w,\mathsf K$ defined in
Section~\ref{sec:coupling}. There is a Borel map
$F:\mathsf W\to\mathsf X$ such that
\[
 \mathsf K(w,\cdot)=\delta_{F(w)}
       \quad\text{for $\mathsf w$-almost every }w,
 \qquad X=F(W)\quad\mathsf Q\text{-almost surely}.
\]
For each fixed $g\in\Aut(T_d)$,
$F(gW)=gF(W)$ almost surely. The free Ising law is the law of a
Borel map of vertex IID uniform labels which commutes with every
automorphism on every label input. The same spin-factor
conclusion holds at $\theta=0$.
\end{proposition}

\begin{proof}
First suppose $\theta>0$. The two uniqueness propositions give
$V(t)=0$ for every deterministic time in their respective parameter
ranges. Invariance permits any vertex in place of the auxiliary
vertex $o$. Intersecting the full-probability events for all vertices
and nonnegative rational times gives
$h_v(q)=h'_v(q)$ simultaneously at those indices. Both field arrays
are continuous, so they agree at every time and vertex. Hence their
magnetizations agree, and the common innovation identity gives
\[
 X_v(t)=W_v(t)+\int_0^t m_v(s)\dd s
       =W_v(t)+\int_0^t m'_v(s)\dd s=X'_v(t).
\]
Thus $\widehat{\mathsf Q}\{X=X'\}=1$ as an equality of full path
configurations. The use of rational times avoided an uncountable
intersection over time.

For $v\in\mathcal V$ and $q\in\mathbb Q_{\ge0}$ choose the finite
Borel version
\[
 f_{v,q}(w)=\int x_v(q)\,\mathsf K(w,\dd x)
\]
when the absolute integral is finite, and set it to zero otherwise.
Kernel integration makes this Borel, and integrability of $X_v(q)$
ensures that the exceptional inputs have $\mathsf w$-measure zero.
Taking expectations in~\eqref{eq:copy-coordinate-variance} and using
$X=X'$ gives
\[
 0=\E_{\widehat{\mathsf Q}}(X_v(q)-X'_v(q))^2
   =2\E_{\mathsf Q}\Var(X_v(q)\mid W).
\]
It follows that $X_v(q)=f_{v,q}(W)$ almost surely. We may impose
these identities simultaneously for all the countably many pairs
$(v,q)$. Disintegrating them shows that, for $\mathsf w$-almost
every $w$, $\mathsf K(w,\cdot)$ gives mass one to
\[
 \{x\in\mathsf X:x_v(q)=f_{v,q}(w)
                 \text{ for every }v,q\}.
\]
This set has at most one element, since continuous paths are
determined by their rational-time values. Its conditional mass one
also makes it nonempty. Thus $\mathsf K(w,\cdot)$ is a point mass
for almost every $w$.

The conditional point-mass principle for measurable recovery also
appears in Kurtz~\cite[Lemma~1.3 in arXiv:1305.6747v2]{Kurtz2014}.
We give a Borel choice of that path, including values on all other
inputs. For each $v$ and integer $j\ge1$, let $p_{v,j}(w)$ be the
piecewise linear path whose value at $k2^{-j}$ is
$f_{v,k2^{-j}}(w)$ for every integer $k\ge0$. These are Borel maps
to $C([0,\infty),\mathbb R)$ with its locally uniform topology.
Let $\mathcal G$ be the set of $w$ for which, at every vertex, these
paths are Cauchy uniformly on every compact interval and
$f_{v,0}(w)=0$. This set is Borel. Indeed, the uniform Cauchy
condition on $[0,N]$ can be tested by the supremum over rational
times there, and the quantifiers over vertices, integers $N$, error
levels, and pairs of sequence indices are all countable. On
$\mathcal G$ the limits are continuous paths starting at zero.
Define $F_v(w)$ to be this limit for $w\in\mathcal G$, and let
$F(w)$ be the all-zero path configuration for $w\notin\mathcal G$.
The limit map on the Cauchy set of a Polish space is Borel, so this
defines a Borel map $F:\mathsf W\to\mathsf X$.

For almost every $w$, the rational coordinates just specified are
those of the unique path configuration charged by $\mathsf K(w,\cdot)$.
Uniform continuity of each such path on compact intervals shows
that its dyadic linear interpolants converge uniformly there.
Consequently $w\in\mathcal G$ and $F(w)$ is that configuration.
This proves both $\mathsf K(w,\cdot)=\delta_{F(w)}$ almost everywhere
and $X=F(W)$ in the original law.

Fix now one automorphism $g$. Invariance of $\mathsf Q$ and
uniqueness of conditional expectations imply
\begin{equation}\label{eq:recovery-coordinate-equivariance}
 f_{v,q}(gw)=f_{g^{-1}v,q}(w)
       \quad\text{for $\mathsf w$-almost every }w
\end{equation}
for each $v,q$. Intersecting over those countably many pairs gives
all these identities simultaneously for this fixed $g$. The dyadic
interpolants then obey
$p_{v,j}(gw)=p_{g^{-1}v,j}(w)$ for every $v,j$. On this full-measure
set, $w\in\mathcal G$ if and only if $gw\in\mathcal G$, because
its defining Cauchy conditions are the same after permuting vertices.
The all-zero default is invariant as well, so
$F(gw)=gF(w)$ almost surely. The full-measure set in this argument
is allowed to depend on $g$; none is asserted for all automorphisms
at once.

Define a spin map on every innovation input by
\begin{equation}\label{eq:spin-factor}
 H_v(w)=
 \begin{cases}
 \displaystyle\lim_{n\to\infty}\frac{F_v(w)(n)}{n},
   &\text{if this limit exists and belongs to $\{-1,1\}$},\\[2mm]
 1,&\text{otherwise}.
 \end{cases}
\end{equation}
Existence of either indicated limit is a Borel condition expressed
through countably many error bounds and integer indices. Thus each
$H_v$ is Borel, and so is $H$ into the countable spin product.
On the planted probability space,
\[
 \frac{F_v(W)(n)}n=\frac{X_v(n)}n
      =\sigma_v+\frac{B_v(n)}n
\]
simultaneously for the integer indices. For every $\varepsilon>0$,
the Gaussian tail bound gives
\[
 \sum_{n\ge1}\P(|B_v(n)|>\varepsilon n)
       \le\sum_{n\ge1}2e^{-\varepsilon^2n/2}<\infty.
\]
Borel--Cantelli, first for positive rational $\varepsilon$ and then
over the countable vertex set, gives $B_v(n)/n\to0$ simultaneously
for every vertex. Therefore $H(W)=\sigma$ almost surely. Under
product Wiener input, $H$ has exactly the free zero-field Ising law.
For the fixed $g$ above, the identity $F(gw)=gF(w)$ makes the
sequences in~\eqref{eq:spin-factor} agree after the same permutation.
Their limits and their failure cases consequently agree. The same
default sign at every site gives $H(gw)=gH(w)$ almost surely.

Finally, Wiener measure is a Borel probability measure on the Polish
space $C_0([0,\infty),\mathbb R)$. The standard Borel sampling
theorem therefore supplies a Borel map
$\xi:[0,1]\to C_0([0,\infty),\mathbb R)$ which sends Lebesgue measure
to Wiener measure. For independent uniform labels $(U_v)$, set
\[
 \Phi(U)=H\bigl((\xi(U_v))_{v\in\mathcal V}\bigr).
\]
Applying the same $\xi$ at every site gives product Wiener law and
commutes pointwise with every vertex permutation. Thus $\Phi$ is
Borel, has law $\mu_{d,\beta}$, and satisfies
$\Phi(gU)=g\Phi(U)$ almost surely for each fixed $g$.
Lemma~\ref{lem:pointwise-spin-factor} replaces $\Phi$ by a Borel map
with the same law that is equivariant on every label input.

When $\theta=0$, the broadcast law is already a product of fair signs.
The rule $\Phi_v(U)=1$ for $U_v\le1/2$ and $\Phi_v(U)=-1$ otherwise
gives that law and commutes with every permutation for every input.
All choices above are for the fixed pair $(d,\beta)$. Together with
the known converse in Section~\ref{sec:introduction}, this proves
Theorem~\ref{thm:main}.
\end{proof}
